% !TeX root = ../教材题库_全部条目_转写工作稿.tex
% Source: XJ-B1-C05-RV-P208-Y1276
% 第5章 / 章末复习 复习题五 / 教材 p.208、209 / PDF 第215、216页
% OCR未校对：公式与文字均可能有误；未完成数学TeX转写；禁止直接发给学生
% 原题及相邻原解截图（核对用）：../题目原图/XJ-B1-C05-RV-P208-Y1276-01.webp, ../题目原图/XJ-B1-C05-RV-P208-Y1276-02.webp
% 本文件只在转写完成后作为题干使用；OCR 条目尚不是可靠题干。
\begin{bankitem}{复习题五 第22题}{第5章 / 章末复习 复习题五 / 教材 p.208、209 / PDF 第215、216页}{XJ-B1-C05-RV-P208-Y1276}{OCR未校对：公式与文字均可能有误；未完成数学TeX转写；禁止直接发给学生}
\begin{ocrdraft}
[教材p.208 / PDF第215页]
\par
22．
\par
已知函数不）=sin()x十@>0，]到
\par
为函数丆0）的零点，
\par
4
\par
为函数）图象的对称轴，且不，）在（一
\par
一）上单调，求的最
\par
5丌
\par
直线一
\par
丌
\par
18，
\par
36
\par
大值．
\par
[教材p.209 / PDF第216页]
\par
上下而求索
\end{ocrdraft}
\end{bankitem}
